\documentclass{homework}
% \usepackage{graphicx} % Required for inserting images
\usepackage{amsmath}
\usepackage{gensymb}
\usepackage{subfig}
\usepackage{float}

\class{ECEN 5478}
\title{Homework 5}
\author{Sean Jennings}
\date{October 26, 2023}

\begin{document}
\maketitle

\section{Problem 1}
If $f$ is a convex function, by the gradient inequality we have
\begin{equation*}
    f(y) \geq f(x) + \nabla f(x) (y - x) \\
\end{equation*}
for all $x,y\in\R^n$, and also
\begin{equation*}
    f(x) \geq f(y) + \nabla f(y) (x - y).
\end{equation*}
Adding together these equations yields
\begin{equation*}
    f(x) + f(y) \geq f(x) + f(y) + \nabla f(x) (y - x) - \nabla f(y) (y - x) \\
\end{equation*}
which reduces to:
\begin{equation*}
    (\nabla f(x) - \nabla f(y)) (x - y) \geq 0
\end{equation*}
showing that $\nabla f$ is a monotone operator.

\section{Problem 2}

If we let $f_i(x) = e^{x_i}$ where $x_i$ is the $i$th entry of $x$, then $f_i$
is convex for each $i$. This is easy to see by noting that the hessian is positive
semi-definite ($\nabla^2 f_i(x) \succeq 0$) for all $x$
since the $i,i$th entry is equal to $e^{x_i} > 0$ and all other entries are 0.

Since each $f_i(x) := e^{x_i}$ is convex, we also have that
$\sum_{i=1}^{n} f_i(x)$ is convex. Since $\log$ is a monotonic function,
then the function of interest is the composition of a monotone operator with
a convex function and is therefore also convex (Beck Thm 7.22).

\section{Problem 3}

The function of interest can be rewritten as:
\begin{equation*}
    f(x) = \max_{1 \leq i \leq n} \{f_i(x)\}
\end{equation*}

where $f_i: \R^n \to \R$ is given by
\begin{equation*}
    f_i(x) = e_i^T x
\end{equation*}
with $e_i$ as the $i$th standard basis vector in $\R^n$.

Each $f_i$ is maps $\R^n$ to a half space and is thus convex (Beck Lemma 6.3), so that $f$ is the
pointwise maximum of convex functions and is therefore also convex (Beck Thm 7.25).

\section{Problem 4}

\begin{letterenum}
\item $Y = \sum_{i=1}^{N} X_i$ is sub-Weibull since it is the sum of subW variables.
    By the closure lemma given in lecture, we have
    \begin{equation*}
        Y \sim \text{subW} \biggl(\max_{1\leq i \leq N} \{\theta_i\},
            \sum_{i=1}^{N} \nu_i\biggr)
    \end{equation*}
    Or equivalently, $\theta_y = \max_{1 \leq i \leq N} \{\theta_i\}$ and
    $\nu_y = \sum_{i=1}^N \nu_i$.

\item In lecture, we also proved that for $a>0$, $\theta > 0$ $\nu > 0$, and $X \sim \subW (\theta, \nu)$
\begin{equation}
    X^a \subW (a \theta, \nu^a \max\{1, a^{a\theta}\})
    \label{eq:subW-x-to-a}
\end{equation}
so that for each $i \in [1, N]$
\begin{equation}
    X_i^2 \sim \subW(2 \theta_i, v_i^2 4^{\theta_i})
    \label{eq:sub-x-squared}
\end{equation}
since $2^{2\theta} = 4^\theta > 4 > 1$ for all $\theta > 0$.

Thus, for $Y = \sum_{i=1}^{N} X_i^2$, we have
\begin{equation}
    Y \sim \subW ( 2 \max_{1\leq i \leq N} \{\theta_i\}, \sum_{i=1}^{N} \nu_i^2 4^{\theta_i})
    \label{eq:sum-x-squared}
\end{equation}
We will call $Y \sim \subW (\theta_b, \nu_b)$ where $\theta_b > 0$ and $\nu_b > 0$ are given by the
above equation, since we will use $\theta_b$ and $\nu_b$ later.

\item We must assume that $X_1$ and $X_2$ are independent. Otherwise, $X_1X_2$ may not be subW.

If they are independent, then by subW closure lemma (d) from lecture and equation \ref{eq:sub-x-squared},  we have

\begin{equation*}
    X_1X_2^2 \sim \subW (\theta_1 + 2 \theta_2, \nu_1 \nu_2^2 4^{\theta_2})
\end{equation*}

and then by scalar multiplication (closure lemma (a)), we have
\begin{equation*}
    Y = \alpha X_1X_2^2 \sim \subW (\theta_1 + 2\theta_2, |\alpha| \nu_1 \nu_2^2 4^{\theta_2})
\end{equation*}
We note that since $\alpha < 0$, we also have $|\alpha| = -\alpha$.

\item $Y = \norm{X}$ where $X = [X_1, X_2, \dots X_N]^T$ can be rewritten as $Y = (\sum_{i=1}^N X_i^2)^{1/2}$.
Since this is just $Z^a$ with $Z$ given by (\ref{eq:sum-x-squared}) we can apply equation (\ref{eq:subW-x-to-a}),
noting this time that since $a = 1/2 < 1$, we have $a^{a\theta} < 1$ for all $\theta > 0$.

\begin{equation*}
    Y \sim \subW \biggl(
        \frac{1}{2} \theta_b,
        \nu_b 
    \biggr)
\end{equation*}
or in terms of $\theta_i$ and $\nu_i$:
\begin{align*}
    Y \sim \subW \biggl(&\max_{1 \leq i \leq N} \{\theta_i\},
        \sum_{i=1}^{N} \nu_i^2 4^{\theta_i}
    \biggr)
\end{align*}
\end{letterenum}

\section{Problem 5}

The optimization problem can be written:
\begin{align*}
    &\min_{u \in \mathcal{U}} \phi(u) + \psi(y) \\
    &\text{subj. to } y = G_u u + G_w w_t \\
    &t \in \mathcal{T}
\end{align*}

where
\begin{align*}
    \mathcal{U} &= [0, u_{max}] \times [0, u_{max}] \subset \R^2 \\
    \mathcal{T} &= \N \cup \{0\} \\
    \phi(u) &= \frac{a_1}{2} \norm{u}^2 \\
    \psi(y) &= \frac{a_2}{2} \norm{y - r}^2.
\end{align*}

The gradient of the two components of the cost function are given by:
\begin{align*}
    \nabla \phi(u) &= a_1 u \\
    \nabla \psi(y) &= a_2 (y - r)
\end{align*}

We can then use the projected gradient method to update the control $u_{t}$
based on the measurement $y_t$. Assuming that $u_0\in\mathcal{U}$ is given
and that the measurement $y_t$ is available for all $t \in \mathcal{T}$,
the algorithm is then:
\begin{align*}
    u_{t+1} &=
        \text{proj}_{\mathcal{U}} \biggl\{u_t - \alpha
        (\nabla \phi(u_t) + G_u^T \nabla \psi(y_t)) \biggr\} \\
        &= \text{proj}_{\mathcal{U}} \biggl\{u_t - \alpha
        (a_1 u_t + a_2 G_u^T (y_t - r)) \biggr\}
\end{align*}

The error of the algorithm is given by the error of the online projected
gradient method:
\begin{equation*}
    \norm{u_t - u_t^*} = \rho \norm{u_0 - u_0^*} + \frac{1}{1 - \rho} \bar{\sigma}
\end{equation*}
where $\bar{\sigma} \in \R$ is the suprenum of the path lengths
\begin{equation*}
    \bar{\sigma} = \sup_{\tau > 0} \norm{u_{\tau + 1} - u_\tau}
\end{equation*}
and $\rho$ is the maximum Lipschitz constant of the algorithmic map. If we define
$f_t: \R^2 \to \R$ to be the total cost $f_t(u) = \phi(u) + \psi(G_u u + G_w w_t)$,
then the algorithmic map is $\mathcal{I} - \nabla f_t$, which has Lipschitz
constant related to the strong convexity and smoothness coefficients of $f_t$.

Since we have,
\begin{align*}
    \nabla f_t(u) &= \nabla \phi(u) + G_u^T \nabla \psi ( G_u u + G_w w_t) \\
    \nabla^2 f_t(u) &= \nabla^2 \phi(u) + G_u^T \nabla^2 \psi ( G_u u + G_w w_t) G_u \\
        &= a_1 I + a_2 G_u^T G_u
\end{align*}
$f_t$ is $\mu$-strongly convex and $L$-smooth with coefficients
\begin{align*}
    L &= \lambda_{max}(a_1 I + a_2 G_u^T G_u) \\
    \mu &= \lambda_{min} (a_1 I + a_2 G_u^T G_u)
\end{align*}
and the algorithmic map has Lipschitz constant
$\rho = \max\{|1 - \alpha L|, |1 - \alpha \mu|\}$, independent of $t$.




\end{document}
